\documentclass[../proyecto.tex]{subfiles}
\graphicspath{{\subfix{../images/}}}
\begin{document}

\chapter{Trabajo avanzado}

% máximo 1 página

\section{Infraestructura comercial}

A nivel individual he realizado la infraestructura comercial y técnica, incluyendo la constitución de mi empresa, Piruetas SpA, que me permite la investigación y desarrollo de los productos asociados a esta tesis doctoral.

\section{Prototipado electrónico}

Me he entrenado en nuevas formas de prototipado electrónico, de automatización de procesos de fabricación de electrónica asistida por máquinas, y de documentación de procesos de diseño electrónico.

\section{Prototipado computacional}

Me he entrenado en nuevas formas de automatización computacional, incluyendo la escritura de código, y la creación automatizada de documentación de código, además de estrategias de publicación de bibliotecas de código.

\section{Cursos de electrónica y computación}

En 2024 creamos en conjunto con el académico y artista Matías Serrano una serie de cursos de pregrado que dictamos en la Escuela de Diseño de la Universidad Diego Portales: el primer semestre dictamos el curso "DIS8644 Taller de diseño de máquinas electrónicas", y en el segundo semestre " DIS8645 Taller de diseño de máquinas computacionales".

La primera versión de cada uno las dictamos en conjunto en el año 2025, y las segundas versiones las estamos dictando este 2026, convirtiéndonos en los talleres de tercer y cuarto año de diseño con mayor lista de espera y cantidad de estudiantes. 

Como parte de nuestras investigaciones doctorales y prácticas conjuntas, empezamos este 2026 a escribir un libro académico de pregrado titulado "Máquinas sonoras", que esperamos sirva de antesala a nuestras tesis doctorales individuales, y que se convierta en un texto de referencia para cursos de enseñanza de sonido, electrónica y programación.

También desde 2026 estoy dictando el curso \textit{DIS09214 Pensamiento computacional}, un nuevo curso obligatorio para estudiantes de segundo año de diseño, que se enfoca en la enseñanza de los paradigmas de la computación, aplicada a la creación de gráficas interactivas con la herramienta p5.js. Mis apuntes del curso también los estoy convirtiendo en un libro académico, que también espero se convierta en un texto de referencia en pregrado en diseño, artes y disciplinas afines.

\end{document}